\chapter{Isolation}
\index{Isolation}

\section*{Definition and Overview}
Isolation, in health care, has two main meanings. Social isolation is the lack of meaningful contact with others, which harms physical and mental health and is increasingly recognised as a major public health problem. Isolation also refers to the public health measure of separating people who have or may have an infectious disease from others, to prevent spread, as used for COVID-19, influenza, and other contagious conditions. This chapter covers both: how to protect wellbeing during illness-related isolation, and how to recognise, prevent, and address social isolation and loneliness, especially in older people, who are most at risk.

\section*{Epidemiology}
Social isolation and loneliness affect around one in five people, and rates are highest among older people, people living alone, carers, people with chronic illness or disability, and those who have recently lost a partner. Loneliness is associated with a 26 per cent increase in the risk of premature death, comparable to the effects of smoking and obesity. During the COVID-19 pandemic, required isolation affected millions and increased loneliness, anxiety, and depression, while demonstrating that illness-related isolation can be managed with planning and connection. Many older people report that loneliness is a regular part of life.

\section*{Aetiology and Pathophysiology}
Social isolation arises from the loss or absence of social connections, which can follow retirement, bereavement, relocation, disability, caring, or life changes, and from barriers to connection such as poor transport, hearing loss, and digital exclusion. Loneliness is the distressing feeling that one's relationships are inadequate, and it is possible to be lonely in a crowd. Both activate stress pathways: chronic loneliness raises cortisol, impairs immunity and sleep, and increases inflammation, contributing to heart disease, stroke, dementia, and depression. Illness-related isolation interrupts daily routine and contact, and can cause boredom, anxiety, and mood change, while protecting others from infection.

\section*{Clinical Features}
\begin{itemize}
    \item Social isolation: little contact with family, friends, or community
    \item Loneliness: feeling disconnected, empty, or left out
    \item Loss of interest, low mood, and anxiety
    \item Poor sleep and reduced appetite
    \item Reduced physical activity and neglect of self-care
    \item Reliance on television, phone, or the internet for company
    \item During illness isolation: boredom, restlessness, and worry about the illness
    \item Physical effects over time: raised blood pressure, weakened immunity, and cognitive decline
\end{itemize}

\section*{Differential Diagnosis}
The effects of isolation can be confused with other conditions.
\begin{itemize}
    \item \textbf{Depression:} low mood, hopelessness, and loss of interest
    \item \textbf{Anxiety:} fear of contact, agoraphobia, or health anxiety
    \item \textbf{Dementia and cognitive decline:} may cause withdrawal
    \item \textbf{Physical illness:} chronic disease causing reduced mobility and contact
    \item \textbf{Hearing and vision loss:} making connection harder
    \item \textbf{Grief and bereavement:} withdrawal after loss
    \item \textbf{Substance use:} isolation associated with dependence
\end{itemize}

\section*{When to Seek Medical Care}
People who feel persistently lonely, low, or anxious, or who are withdrawing from life, should seek help from a GP, psychologist, or community service. Those in illness-related isolation should know the signs of worsening illness and when to seek medical care. Family and friends should check in on isolated people, especially older relatives and neighbours.

\textbf{Contact your local emergency services for:}
\begin{itemize}
    \item Thoughts of suicide or self-harm
    \item Difficulty breathing or severe illness during isolation
    \item Confusion, collapse, or inability to care for oneself
    \item Any situation in which a person is in immediate danger
\end{itemize}

\section*{Diagnosis and Investigations}
\begin{itemize}
    \item \textbf{Screening:} simple questions about loneliness, social contact, and support
    \item \textbf{Assessment of contributing factors:} hearing, vision, mobility, transport, and digital access
    \item \textbf{Mental health assessment:} depression, anxiety, and grief
    \item \textbf{Review of physical health:} conditions and medicines that affect function
    \item \textbf{Social assessment:} housing, finances, family, and community connections
    \item \textbf{Risk assessment:} safety and self-care in people living alone
\end{itemize}

\section*{Management}
\begin{itemize}
    \item \textbf{Connection strategies:} regular phone calls, visits, and community activities
    \item \textbf{Community programs:} social groups, exercise classes, and volunteering
    \item \textbf{Digital connection:} video calls, online groups, and digital skills support
    \item \textbf{Treating barriers:} hearing aids, mobility aids, transport, and home visits
    \item \textbf{Psychological support:} counselling and cognitive behavioural therapy for loneliness and mood
    \item \textbf{Peer support:} befriending services and support lines
    \item \textbf{Planning for illness isolation:} supplies, medicines, a support contact, and a plan for worsening symptoms
    \item \textbf{Maintaining routine:} structure, activity, and time outdoors during isolation
\end{itemize}

\section*{Complications}
\begin{itemize}
    \item Depression, anxiety, and suicide risk
    \item Heart disease, stroke, and dementia
    \item Weakened immunity and slower recovery from illness
    \item Falls and injuries in older people without support
    \item Neglect of self-care, medicines, and appointments
    \item Worsening of chronic disease
    \item Deterioration of physical and mental health during prolonged illness isolation
\end{itemize}

\section*{Prognosis and Outcomes}
Social isolation is not a fixed state, and loneliness responds well to connection. Most people who engage in community activities, maintain regular contact, and address barriers such as hearing or mobility report substantial improvement in mood and wellbeing. During illness-related isolation, most people cope well with planning and support, and symptoms of anxiety and loneliness settle as they recover. The key is early recognition and action, both for individuals and for the friends, families, and services around them.

\section*{Prevention and Self-Care}
\begin{itemize}
    \item Keep regular contact with family, friends, and neighbours
    \item Join community, faith, or interest groups
    \item Volunteer or take part in local activities
    \item Learn video calling and stay in touch digitally if that suits you
    \item Address hearing, vision, and mobility problems
    \item If isolating due to illness, plan ahead: supplies, medicines, and a support person
    \item Keep a routine of meals, activity, and sleep during isolation
    \item Reach out early if you feel lonely or low
    \item Check on isolated neighbours and relatives
    \item Seek support from befriending and community services
\end{itemize}

\section*{Sources and Further Reading}
The following reputable sources provide further information for healthcare students and patients seeking reliable, current advice about isolation and loneliness.
\begin{itemize}
    \item Australian Institute of Health and Welfare. \textit{Social isolation and loneliness.} https://www.aihw.gov.au/
    \item Beyond Blue. \textit{Support for loneliness and low mood.} https://www.beyondblue.org.au/
    \item EveryAGE Counts. \textit{Ageism and loneliness campaign.} https://www.everyagecounts.org.au/
    \item Healthdirect Australia. \textit{Loneliness and social isolation.} https://www.healthdirect.gov.au/
    \item Relationships Australia. \textit{Connection and support services.} https://www.relationships.org.au/
    \item Campaign to End Loneliness. \textit{Research and resources.} https://www.campaigntoendloneliness.org/
\end{itemize}
